% ECONOMICS_PROBLEM: {"id": "C78-619", "title": "A multiplayer extension of Weber-law bargaining", "jel_code": "C78", "source_id": "OP-0223"}
% !TeX program = xelatex
\documentclass[11pt,a4paper]{article}
\usepackage[margin=23mm]{geometry}
\usepackage{fontspec}
\setmainfont{Latin Modern Roman}
\usepackage{amsmath,amssymb,mathtools}
\usepackage{xcolor,enumitem,booktabs,longtable,array,xurl,hyperref}
\definecolor{muted}{HTML}{53636C}
\hypersetup{colorlinks=true,urlcolor=blue,linkcolor=blue}
\setlength{\parindent}{0pt}
\setlength{\parskip}{5pt}
\newcommand{\R}{\mathbb R}
\newcommand{\E}{\mathbb E}
\newcommand{\N}{\mathbb N}
\newcommand{\ind}{\mathbf 1}
\newcommand{\Var}{\operatorname{Var}}
\newcommand{\Cov}{\operatorname{Cov}}
\newcommand{\supp}{\operatorname{supp}}
\DeclareMathOperator*{\argmin}{arg\,min}
\DeclareMathOperator*{\argmax}{arg\,max}
\newcommand{\entrymeta}[3]{{\sffamily\small\color{muted}\textbf{\texttt{#1}}\quad|\quad #2\par #3\par}}
\newcommand{\Problemref}[1]{\texttt{#1}}
\newcommand{\JELref}[2]{\texttt{#2} (JEL \texttt{#1})}
\begin{document}
\section*{A multiplayer extension of Weber-law bargaining}
\textbf{Problem ID:} \texttt{C78-619}\quad \textbf{JEL:} \texttt{C78}\par
% BEGIN ECONOMICS STATEMENT
\label{problem:OP-0223}
% GAME_THEORY_PROBLEM: {"local_id": "GT-093", "original_number": 93, "kind": "R", "entry_kind": "statement_only", "published": false, "review_required": true, "category": "game_theory"}
\label{game:GT-093}
{\sffamily\small\color{muted}
\noindent\textbf{ID:} GT-093. \textbf{Category:} Bargaining and behavioural game theory.
\textbf{Status:} R; the source states a research direction rather than a canonical conjecture.
\par}
\subsection*{Definitions and mathematical statement}
The following is an explicitly editorial benchmark. Fix $n\geq3$, a resource $X>0$, and continuously differentiable, strictly increasing concave utilities $U_i:[0,X]\to\mathbb R$. Put $d_i=U_i(0)$ and $g_i(t)=U_i(t)-d_i$. The efficient resource allocations form
\[
 \mathcal A_X=\{x\in\mathbb R_{\geq0}^n:\sum_i x_i=X\}.
\]
For constants $0<\kappa_i<1$ and a positive baseline gain $b_i$, define the Weber acceptance floor and acceptance set by
\[
 a_i(b_i)=(1-\kappa_i)b_i,\qquad
 \mathcal A_X(b)=\{x\in\mathcal A_X:g_i(x_i)\geq a_i(b_i)\text{ for all }i\}.
\]
This uses utility gains from disagreement; the corresponding just-noticeable loss is $\kappa_i b_i$.

A mathematical extension should specify a proposal-and-baseline update protocol, its admissible histories, and a selection map
$\Phi(U,d,\kappa,X)\in\mathcal A_X$, then prove convergence of the protocol and characterize its selected allocation. For a definite benchmark, require permutation covariance and invariance of the resource allocation under $U_i\mapsto\alpha_iU_i+\beta_i$, $d_i\mapsto\alpha_id_i+\beta_i$, $\alpha_i>0$. Require also the following susceptibility limit: for every $\theta_i>0$, with $\kappa_i=\varepsilon\theta_i$ and $\varepsilon\downarrow0$,
\[
 \Phi(U,d,\varepsilon\theta,X)\longrightarrow
 \operatorname*{arg\,max}_{x\in\mathcal A_X,\ x_i>0}
          \sum_{i=1}^n w_i\log g_i(x_i),\qquad
 w_i=\frac{\theta_i^{-1}}{\sum_j\theta_j^{-1}}.
\]
The displayed maximizer is unique under these utility assumptions. The target is to construct and analyze a protocol using the stated acceptance floors, including how baselines change, or to identify incompatible desired properties.

\subsection*{Short English statement}
Specify how three or more players exchange and revise Weber-based offers, and prove what allocation their negotiations select.

\subsection*{Variants and scope boundaries}
Weber's law constrains discrimination or tolerated changes; it does not by itself determine a unique bargaining outcome or a baseline-update rule. Defining a weighted Nash maximizer alone does not supply the requested behavioural protocol. The inverse-$\theta_i$ weights extend the bilateral susceptibility weights, but their multiplayer use here is an editorial requirement, not a source theorem. Nonconcave utilities, lotteries, and probabilistic discrimination are separate models.

\subsection*{Source relationship and status}
The source gives the discrimination threshold in Section 1.2, affine-covariant utility-gain acceptance in Section 3, bilateral proposal dynamics and partial allocations in Sections 4--5, and a multiplayer direction in Section 7. It does not supply a unique $n$-player protocol or the displayed $n$-player limit. Accordingly this entry states a concrete benchmark and does not label its chosen axioms as an authorial open conjecture.

\subsection*{References}
\begingroup\footnotesize\raggedright
\noindent [S1] V.~G. Bardakhchyan and A.~E. Allahverdyan, \emph{Bargaining via Weber's law} (2024; inspected revision v3, 2025), arXiv:2408.02492v3. Inspected locators: Section 1.2, Equation (1); Section 3, Equations (9)--(16); Sections 4--5; Section 7.
\url{https://arxiv.org/html/2408.02492v3}
\par\endgroup

\subsection*{Common conventions: Source mathematical conventions}

\textbf{Finite games.} For a finite set $A$, $\Delta(A)$ is its probability
simplex. Players choose independent mixed strategies in Nash equilibrium;
correlation is allowed only when explicitly part of the solution concept.
In a game with payoff functions $u_i$ and strategy profile $x$, an additive
$\varepsilon$ Nash equilibrium satisfies
\[
 u_i(x)\geq u_i(x_i',x_{-i})-\varepsilon
 \quad\text{for every player $i$ and every admissible deviation $x_i'$}.
\]
For numerical approximation, payoffs are normalized as locally specified.
An $\varepsilon$ payoff residual is distinct from distance at most
$\varepsilon$ to the set of exact equilibria. Point convergence, convergence
of empirical play, and convergence in distance to an equilibrium set are
also distinct.

\textbf{Computational representations.} Unless stated otherwise, a complexity
question takes rational data with binary encoding; polynomial time is measured
in total bit length. A fixed-accuracy polynomial algorithm is not necessarily
polynomial in $1/\varepsilon$, and neither is necessarily polynomial in
$\log(1/\varepsilon)$. Succinct games and value, demand, or sampling oracles
must declare their interfaces. Exact real solutions require an explicit
representation; an unspecified list of arbitrary real numbers is not a
finite computational output. A claimed algorithmic barrier is meaningful
only for its specified model and reductions.

\textbf{Dynamic games.} Histories, information, observation, and allowable
randomization are part of the model. A stationary strategy depends only on
the current state; a positional strategy in a deterministic graph game
chooses one action at each controlled vertex. Nash equilibrium at an initial
state and equilibrium after every history are different requirements.
An expectation of a pathwise $\liminf$ payoff, a $\liminf$ of expected averages,
a limiting expected average, and a uniform finite-horizon equilibrium are
different criteria. None is silently substituted for another.

\textbf{Allocation and mechanisms.} A complete allocation assigns every item
exactly once unless otherwise stated. Zero-valued goods, negative values,
capacity constraints, feasibility, lotteries, and copies of goods affect
fairness definitions and are specified locally. Dominant-strategy incentive
compatibility, Bayesian incentive compatibility, universal truthfulness,
and truthfulness in expectation impose different inequalities. Whenever
an objective ratio has zero denominator, its convention is stated or those
instances are excluded.

\textbf{Infinite-dimensional models.} Expectations, integrals, stochastic
equations, and optimization objectives require their local regularity,
integrability, filtrations, and admissible domains. A backward--forward
mean-field system is a boundary-value problem; it is not an initial-value
semiflow without an additional construction. Long-time, large-population,
and vanishing-noise limits require an order of limits and a convergence
topology. If a minimum need not be attained, the objective is its infimum
or supremum and the approximation requirement is stated separately.
% END ECONOMICS STATEMENT
\end{document}
